% Options for packages loaded elsewhere
% Options for packages loaded elsewhere
\PassOptionsToPackage{unicode}{hyperref}
\PassOptionsToPackage{hyphens}{url}
\PassOptionsToPackage{dvipsnames,svgnames,x11names}{xcolor}
%
\documentclass[
  letterpaper,
  DIV=11,
  numbers=noendperiod]{scrartcl}
\usepackage{xcolor}
\usepackage{amsmath,amssymb}
\setcounter{secnumdepth}{-\maxdimen} % remove section numbering
\usepackage{iftex}
\ifPDFTeX
  \usepackage[T1]{fontenc}
  \usepackage[utf8]{inputenc}
  \usepackage{textcomp} % provide euro and other symbols
\else % if luatex or xetex
  \usepackage{unicode-math} % this also loads fontspec
  \defaultfontfeatures{Scale=MatchLowercase}
  \defaultfontfeatures[\rmfamily]{Ligatures=TeX,Scale=1}
\fi
\usepackage{lmodern}
\ifPDFTeX\else
  % xetex/luatex font selection
\fi
% Use upquote if available, for straight quotes in verbatim environments
\IfFileExists{upquote.sty}{\usepackage{upquote}}{}
\IfFileExists{microtype.sty}{% use microtype if available
  \usepackage[]{microtype}
  \UseMicrotypeSet[protrusion]{basicmath} % disable protrusion for tt fonts
}{}
\makeatletter
\@ifundefined{KOMAClassName}{% if non-KOMA class
  \IfFileExists{parskip.sty}{%
    \usepackage{parskip}
  }{% else
    \setlength{\parindent}{0pt}
    \setlength{\parskip}{6pt plus 2pt minus 1pt}}
}{% if KOMA class
  \KOMAoptions{parskip=half}}
\makeatother
% Make \paragraph and \subparagraph free-standing
\makeatletter
\ifx\paragraph\undefined\else
  \let\oldparagraph\paragraph
  \renewcommand{\paragraph}{
    \@ifstar
      \xxxParagraphStar
      \xxxParagraphNoStar
  }
  \newcommand{\xxxParagraphStar}[1]{\oldparagraph*{#1}\mbox{}}
  \newcommand{\xxxParagraphNoStar}[1]{\oldparagraph{#1}\mbox{}}
\fi
\ifx\subparagraph\undefined\else
  \let\oldsubparagraph\subparagraph
  \renewcommand{\subparagraph}{
    \@ifstar
      \xxxSubParagraphStar
      \xxxSubParagraphNoStar
  }
  \newcommand{\xxxSubParagraphStar}[1]{\oldsubparagraph*{#1}\mbox{}}
  \newcommand{\xxxSubParagraphNoStar}[1]{\oldsubparagraph{#1}\mbox{}}
\fi
\makeatother


\usepackage{longtable,booktabs,array}
\newcounter{none} % for unnumbered tables
\usepackage{calc} % for calculating minipage widths
% Correct order of tables after \paragraph or \subparagraph
\usepackage{etoolbox}
\makeatletter
\patchcmd\longtable{\par}{\if@noskipsec\mbox{}\fi\par}{}{}
\makeatother
% Allow footnotes in longtable head/foot
\IfFileExists{footnotehyper.sty}{\usepackage{footnotehyper}}{\usepackage{footnote}}
\makesavenoteenv{longtable}
\usepackage{graphicx}
\makeatletter
\newsavebox\pandoc@box
\newcommand*\pandocbounded[1]{% scales image to fit in text height/width
  \sbox\pandoc@box{#1}%
  \Gscale@div\@tempa{\textheight}{\dimexpr\ht\pandoc@box+\dp\pandoc@box\relax}%
  \Gscale@div\@tempb{\linewidth}{\wd\pandoc@box}%
  \ifdim\@tempb\p@<\@tempa\p@\let\@tempa\@tempb\fi% select the smaller of both
  \ifdim\@tempa\p@<\p@\scalebox{\@tempa}{\usebox\pandoc@box}%
  \else\usebox{\pandoc@box}%
  \fi%
}
% Set default figure placement to htbp
\def\fps@figure{htbp}
\makeatother





\setlength{\emergencystretch}{3em} % prevent overfull lines

\providecommand{\tightlist}{%
  \setlength{\itemsep}{0pt}\setlength{\parskip}{0pt}}



 


\KOMAoption{captions}{tableheading}
\makeatletter
\@ifpackageloaded{caption}{}{\usepackage{caption}}
\AtBeginDocument{%
\ifdefined\contentsname
  \renewcommand*\contentsname{Table of contents}
\else
  \newcommand\contentsname{Table of contents}
\fi
\ifdefined\listfigurename
  \renewcommand*\listfigurename{List of Figures}
\else
  \newcommand\listfigurename{List of Figures}
\fi
\ifdefined\listtablename
  \renewcommand*\listtablename{List of Tables}
\else
  \newcommand\listtablename{List of Tables}
\fi
\ifdefined\figurename
  \renewcommand*\figurename{Figure}
\else
  \newcommand\figurename{Figure}
\fi
\ifdefined\tablename
  \renewcommand*\tablename{Table}
\else
  \newcommand\tablename{Table}
\fi
}
\@ifpackageloaded{float}{}{\usepackage{float}}
\floatstyle{ruled}
\@ifundefined{c@chapter}{\newfloat{codelisting}{h}{lop}}{\newfloat{codelisting}{h}{lop}[chapter]}
\floatname{codelisting}{Listing}
\newcommand*\listoflistings{\listof{codelisting}{List of Listings}}
\makeatother
\makeatletter
\makeatother
\makeatletter
\@ifpackageloaded{caption}{}{\usepackage{caption}}
\@ifpackageloaded{subcaption}{}{\usepackage{subcaption}}
\makeatother
\usepackage{bookmark}
\IfFileExists{xurl.sty}{\usepackage{xurl}}{} % add URL line breaks if available
\urlstyle{same}
\hypersetup{
  pdftitle={Chess-Based Training and Adaptive Functioning: A Study with Italian Preschool and Primary School Children},
  pdfauthor={Nome Autore},
  colorlinks=true,
  linkcolor={blue},
  filecolor={Maroon},
  citecolor={Blue},
  urlcolor={Blue},
  pdfcreator={LaTeX via pandoc}}


\title{Chess-Based Training and Adaptive Functioning: A Study with
Italian Preschool and Primary School Children}
\author{Nome Autore}
\date{2026-03-12}
\begin{document}
\maketitle


\section{Abstract}\label{abstract}

This study examined whether an eight-session Giant Chess Game training
protocol was associated with changes in cognitive, navigational, and
adaptive-functioning outcomes in Italian children aged 5 to 7 years. The
final analytic database included 101 children (57 experimental, 44
control). Minefield route-efficiency scores were recalculated from raw
data when available and otherwise retained from the historical database.
Analyses were rerun with the final database using a reproducible Python
pipeline. PRE scores were first screened for sex and age differences;
training effects were then tested as between-group differences in
PRE-to-POST change with HC3 robust standard errors, with age-adjusted
ANCOVA models as sensitivity analyses. No group x time effect survived
FDR correction across the 45 PRE-POST outcomes. Baseline performance
showed expected age-related differences, particularly for Raven
accuracy, Tower of London rule violations, and Moles Task working-memory
measures. Predictor analyses suggested that teacher-rated adaptive
functioning was associated with selected cognitive-effectiveness
outcomes, most clearly for ABAS-II Health and Safety predicting Tower of
London accuracy effectiveness. Overall, the final analyses did not
provide corrected evidence for a generalized training effect, but they
documented meaningful developmental variation in baseline navigational
and executive-function measures.

\section{Introduzione}\label{introduzione}

Chess-based educational activities have been proposed as a promising
context for training planning, working memory, inhibitory control, and
adaptive problem solving in childhood. The present study tested a
school-based Giant Chess Game protocol designed for preschool and early
primary-school children. The intervention used a large physical
chessboard to make spatial relations, piece movements, and rule-based
planning concrete and embodied. The study focused on whether children
who received the chess protocol showed greater PRE-to-POST improvement
than children in a matched control condition, and whether
adaptive-functioning and everyday executive-function ratings were
associated with individual cognitive change.

\section{Methods}\label{methods}

\subsection{Sample}\label{sample}

The source recruitment sample comprised 104 typically developing Italian
children aged between 5 and 7 years. After database cleaning and
availability checks, the final analytic database included 101 children
(55 boys, 46 girls; mean age = 6.43 years, SD = 0.65). Fifty-seven
children were assigned to the Experimental Chess Protocol Group (GSPS),
and 44 were assigned to the Control Group (CG). The final age
distribution was 9 children aged 5 years, 40 aged 6 years, and 52 aged 7
years.

Participants were recruited from the third year of two preschools (``D.
Ricci'' in Acquasparta and ``G.S. Bernardo'' in Narni) and from the
first and second years of four primary schools (``A.B. Sabin'' in
Acquasparta, Montecchio Primary School, ``Narni S. Lucia'', and ``Narni
Centro''). To support equivalence between groups, each experimental
class was matched with a control class of the same school level. After
testing, children in the control group also participated in chess
activities. Inclusion criteria were: native Italian speakers, no primary
sensory deficits, no neurological conditions, and no significant
psychiatric or developmental disorders. General cognitive functioning
was assessed with Raven's Coloured Progressive Matrices, a non-verbal
measure of intellectual ability (@raven1995). A Raven cut-off
sensitivity check was also conducted using age-specific fifth-percentile
thresholds. None of the children had prior structured chess experience.
The study adhered to the ethical principles of the Declaration of
Helsinki and was approved by the local ethics committee. Written
informed consent was obtained from parents, and assent from each child.

\subsection{Testing at T0 and T1}\label{testing-at-t0-and-t1}

Children were individually tested in a quiet room at school. The GSPS
group was assessed immediately before (T0) and after (T1) participating
in the Giant Chess Game training. The control group, which continued
regular classroom activities without chess practice, was also tested at
T0 and T1. The order of tests was counterbalanced for each participant.

\subsubsection{Corsi Block-Tapping Test}\label{corsi-block-tapping-test}

The Corsi Block-Tapping Test (CBT; @monaco2013span; @corsi1972) was used
to assess visuospatial working memory and short-term memory. The test
uses a board with nine wooden blocks arranged in a scattered pattern.
During the task, the examiner taps a sequence of blocks at a fixed pace,
and the child reproduces the same sequence. The test begins with
sequences of two blocks and increases progressively. The span score
corresponds to the longest sequence level correctly reproduced according
to the task rules. Parallel versions were used at T0 and T1.

\subsubsection{Tower of London}\label{tower-of-london}

The Tower of London (TOL; Shallice, 1982), included in the BVN 5-11
battery (Bisiacchi et al., 2005), assesses planning, problem solving,
and the ability to anticipate consequences. Children must reproduce
target configurations using colored balls and pegs, while respecting
task rules about move order, peg capacity, and minimal solutions.

\subsubsection{Moles Task}\label{moles-task}

The Moles Task is a computerized child-friendly task adapted with moles
as obstacles. It measures spatial memory, navigation skills, and
visuospatial planning strategies on an 8 x 8 grid. Three conditions were
administered:

\begin{itemize}
\tightlist
\item
  In the working-memory condition, children viewed mole positions,
  retained them, and then recalled the exact locations.
\item
  In the planning condition, children traced the shortest path from
  start to finish while avoiding visible obstacles.
\item
  In the working-memory-planning condition, obstacles disappeared before
  response, requiring children to remember their positions and plan an
  efficient route.
\end{itemize}

Raw Minefield data were rescored when available. Minimum paths were
recalculated from \texttt{configuration\_trial.xlsx}, allowing
orthogonal movement only and treating obstacle cells as blocked. The
current paper analyses used
\texttt{data/FINAL\_DATABASE\_RECALCULATED\_MINEFIELD.xlsx}, which
preserves the original database schema and updates only Planning and
working-memory-planning route-efficiency percentages when raw
recalculation is available.

\subsubsection{Teacher Questionnaires}\label{teacher-questionnaires}

The ABAS-II (Adaptive Behavior Assessment System, Second Edition;
@harrison2003; @ferri2014) was used to assess adaptive functioning in
the school context. The teacher form provides scores for adaptive
domains and broader composite indices. The BRIEF-2 (Behavior Rating
Inventory of Executive Function, Second Edition; @gioia2015) was used to
assess everyday executive functioning. Items are rated by teachers and
provide clinical scales and composite indices covering behavioral,
emotional, and cognitive regulation.

\subsection{Giant Chess Game Training}\label{giant-chess-game-training}

The Giant Chess Game consisted of 8 weekly chess lessons, each lasting 1
hour and divided into three work units. Lessons were conducted by a
chess master affiliated with FIDE. The large chessboard measured 2.5 m x
2.5 m and was divided into 8 rows and 8 columns.

\subsubsection{Work Unit 1}\label{work-unit-1}

Children first explored the giant chessboard freely and were guided
through simple spatial movements, such as one step forward, two steps
backward, or one step to the right. The teacher introduced squares,
rows, and diagonals through embodied play activities.

\subsubsection{Work Unit 2}\label{work-unit-2}

The teacher introduced the King and Queen. Children practiced movement
rules through structured games such as ``Catch the King'', and through
problem situations requiring them to reason about controlled space on
the board.

\subsubsection{Work Unit 3}\label{work-unit-3}

The Knight was introduced through movement games emphasizing its
L-shaped displacement. Children practiced rotating the movement pattern
in different spatial directions and adapting it to the board boundaries.

\section{Results}\label{results}

Analyses were rerun using the final paper database,
\texttt{data/FINAL\_DATABASE\_RECALCULATED\_MINEFIELD.xlsx}, through the
reproducible \texttt{analysis\_2} pipeline. The final analysis dataset
included 101 participants. PRE scores were screened for demographic
differences using Welch independent-samples t-tests. Sex differences
were tested by comparing boys and girls on each PRE outcome. Age
differences were tested with pairwise comparisons among 5-, 6-, and
7-year-old children. Multiple comparisons were corrected using the
Benjamini-Hochberg false-discovery-rate (FDR) procedure.

Training effects were tested across 45 PRE-POST outcomes. The primary
model compared PRE-to-POST change between groups using HC3 robust
standard errors, which is equivalent to the group x time interaction
when there are two occasions. Age-adjusted ANCOVA sensitivity models
tested POST scores as a function of PRE score, group, and age. Tables
report uncorrected p values as p raw, together with FDR-adjusted q
values and effect-size estimates when available.

As a preliminary data-quality check, Raven PRE accuracy was screened
against age-specific fifth-percentile cut-offs. Participants were
flagged when Raven PRE accuracy was \textless= 10 at age 5, \textless=
11 at age 6, or \textless= 12 at age 7. Two participants in the final
analytic database fell below these cut-offs: one 5-year-old child in the
experimental group and one 7-year-old child in the control group. This
check was retained as a sensitivity/data-quality audit and was not used
to change the primary paper database.

Time-based outcomes were logarithmically transformed before analysis
when appropriate.

\subsection{Training's Effect on Cognitive
Abilities}\label{trainings-effect-on-cognitive-abilities}

Before testing training effects, PRE scores were screened for possible
demographic differences. No sex-related difference survived FDR
correction. Age-related differences survived FDR correction for eight
PRE outcomes: raw and log-transformed Planning execution/total time,
Raven accuracy, Tower of London rule violations, Moles Task
working-memory accuracy, Moles Task working-memory span, and Moles Task
working-memory trials. The significant age comparisons are summarized
below.

{\def\LTcaptype{none} % do not increment counter
\begin{longtable}[]{@{}
  >{\raggedright\arraybackslash}p{(\linewidth - 12\tabcolsep) * \real{0.1154}}
  >{\raggedright\arraybackslash}p{(\linewidth - 12\tabcolsep) * \real{0.1154}}
  >{\raggedleft\arraybackslash}p{(\linewidth - 12\tabcolsep) * \real{0.1538}}
  >{\raggedleft\arraybackslash}p{(\linewidth - 12\tabcolsep) * \real{0.1538}}
  >{\raggedleft\arraybackslash}p{(\linewidth - 12\tabcolsep) * \real{0.1538}}
  >{\raggedleft\arraybackslash}p{(\linewidth - 12\tabcolsep) * \real{0.1538}}
  >{\raggedleft\arraybackslash}p{(\linewidth - 12\tabcolsep) * \real{0.1538}}@{}}
\toprule\noalign{}
\begin{minipage}[b]{\linewidth}\raggedright
PRE outcome
\end{minipage} & \begin{minipage}[b]{\linewidth}\raggedright
Significant age comparison
\end{minipage} & \begin{minipage}[b]{\linewidth}\raggedleft
Younger group M
\end{minipage} & \begin{minipage}[b]{\linewidth}\raggedleft
Older group M
\end{minipage} & \begin{minipage}[b]{\linewidth}\raggedleft
Hedges g
\end{minipage} & \begin{minipage}[b]{\linewidth}\raggedleft
p raw
\end{minipage} & \begin{minipage}[b]{\linewidth}\raggedleft
FDR q
\end{minipage} \\
\midrule\noalign{}
\endhead
\bottomrule\noalign{}
\endlastfoot
PLANNING\_TR\_LN & 7 vs.~6 & 5.80 & 5.27 & -1.13 & \textless{} .001 &
\textless{} .001 \\
PLANNING\_TE\_LN & 7 vs.~6 & 5.57 & 5.03 & -1.11 & \textless{} .001 &
\textless{} .001 \\
PLANNING\_TR & 7 vs.~6 & 359.41 & 216.68 & -1.14 & \textless{} .001 &
\textless{} .001 \\
Raven\_ACC & 7 vs.~6 & 19.00 & 22.83 & .83 & \textless{} .001 & .005 \\
WM\_TRIAL & 7 vs.~5 & 3.00 & 4.54 & 1.11 & \textless{} .001 & .011 \\
TOL\_VIO\_REG & 7 vs.~5 & 5.00 & 2.38 & -1.21 & \textless{} .001 &
.011 \\
WM\_SPAN & 7 vs.~5 & 1.89 & 2.94 & .97 & \textless{} .001 & .011 \\
WM\_TRIAL & 7 vs.~6 & 3.51 & 4.54 & .72 & .001 & .020 \\
WM\_SPAN & 7 vs.~6 & 2.14 & 2.94 & .71 & .001 & .021 \\
WM\_ACC & 7 vs.~6 & 1.35 & 2.44 & .69 & .002 & .032 \\
\end{longtable}
}

Across the 45 PRE-POST outcomes tested, no group x time effect survived
FDR correction. The strongest uncorrected effects are reported below.

{\def\LTcaptype{none} % do not increment counter
\begin{longtable}[]{@{}lrlrrr@{}}
\toprule\noalign{}
Outcome & Group x time beta & 95\% CI & partial eta2 & p raw & FDR q \\
\midrule\noalign{}
\endhead
\bottomrule\noalign{}
\endlastfoot
CBT\_F\_SPAN & -.64 & {[}-1.16, -.12{]} & .056 & .015 & .452 \\
WM\_ACC & .93 & {[}.15, 1.72{]} & .054 & .020 & .452 \\
TOL\_TE & -18.38 & {[}-36.29, -.47{]} & .040 & .044 & .665 \\
TOL\_TE\_LN & -.13 & {[}-.26, .01{]} & .034 & .064 & .716 \\
PWM\_EFF & -2.66 & {[}-6.35, 1.02{]} & .020 & .157 & .811 \\
PLANNING\_MATTI\_MIN & -3.23 & {[}-7.94, 1.49{]} & .019 & .179 & .811 \\
PLANNING\_TP\_LN & .25 & {[}-.12, .62{]} & .018 & .191 & .811 \\
PWM\_PERC & 7.03 & {[}-3.52, 17.57{]} & .017 & .192 & .811 \\
\end{longtable}
}

The age-adjusted ANCOVA sensitivity analyses also found no group effect
surviving FDR correction. Thus, the final analysis did not provide
corrected evidence that the experimental group improved more than the
control group from T0 to T1.

\subsection{Relationship between Autonomy and Executive
Functions}\label{relationship-between-autonomy-and-executive-functions}

To examine whether adaptive functioning and everyday executive
functioning were associated with cognitive change, ABAS-II and BRIEF-2
scores were tested as predictors of individual treatment effectiveness
across bounded cognitive outcomes. Predictor models were adjusted for
age, and p values were corrected using the Benjamini-Hochberg FDR
procedure.

One adaptive-functioning predictor survived FDR correction within the
ABAS-II/BRIEF-2 predictor set: higher ABAS-II supplemental Health and
Safety scores predicted greater Tower of London accuracy effectiveness.
Additional nominal associations are shown below and should be
interpreted cautiously when the FDR-corrected q value is above .05.

{\def\LTcaptype{none} % do not increment counter
\begin{longtable}[]{@{}
  >{\raggedright\arraybackslash}p{(\linewidth - 12\tabcolsep) * \real{0.1200}}
  >{\raggedright\arraybackslash}p{(\linewidth - 12\tabcolsep) * \real{0.1200}}
  >{\raggedleft\arraybackslash}p{(\linewidth - 12\tabcolsep) * \real{0.1600}}
  >{\raggedright\arraybackslash}p{(\linewidth - 12\tabcolsep) * \real{0.1200}}
  >{\raggedleft\arraybackslash}p{(\linewidth - 12\tabcolsep) * \real{0.1600}}
  >{\raggedleft\arraybackslash}p{(\linewidth - 12\tabcolsep) * \real{0.1600}}
  >{\raggedleft\arraybackslash}p{(\linewidth - 12\tabcolsep) * \real{0.1600}}@{}}
\toprule\noalign{}
\begin{minipage}[b]{\linewidth}\raggedright
Predictor
\end{minipage} & \begin{minipage}[b]{\linewidth}\raggedright
Outcome
\end{minipage} & \begin{minipage}[b]{\linewidth}\raggedleft
beta
\end{minipage} & \begin{minipage}[b]{\linewidth}\raggedright
95\% CI
\end{minipage} & \begin{minipage}[b]{\linewidth}\raggedleft
p raw
\end{minipage} & \begin{minipage}[b]{\linewidth}\raggedleft
FDR q
\end{minipage} & \begin{minipage}[b]{\linewidth}\raggedleft
adj. R2
\end{minipage} \\
\midrule\noalign{}
\endhead
\bottomrule\noalign{}
\endlastfoot
ABAS\_Health\_saf\_suppongo & Tower of London accuracy effectiveness &
15.30 & {[}7.83, 22.77{]} & \textless{} .001 & .005 & .220 \\
BRIEF\_Em\_Con & Planning-working-memory route-efficiency effectiveness
& 28.35 & {[}9.54, 47.15{]} & .003 & .105 & .061 \\
ABAS\_school\_liv\_suppongo & Tower of London accuracy effectiveness &
15.72 & {[}4.93, 26.51{]} & .004 & .136 & .148 \\
ABAS\_Health\_saf\_sic & Raven accuracy effectiveness & -7.49 &
{[}-12.79, -2.18{]} & .006 & .165 & .041 \\
ABAS\_fun\_acc & Planning composite effectiveness & 43.16 & {[}9.89,
76.43{]} & .011 & .241 & .080 \\
ABAS\_comm\_use & Planning composite effectiveness & 31.44 & {[}5.59,
57.29{]} & .017 & .320 & .040 \\
ABAS\_Leisure & Planning-working-memory composite effectiveness & 10.42
& {[}1.28, 19.56{]} & .025 & .409 & .066 \\
BRIEF\_ERI & Planning-working-memory route-efficiency effectiveness &
20.24 & {[}2.46, 38.03{]} & .026 & .409 & .040 \\
ABAS\_Leisure & Planning-working-memory accuracy effectiveness & 11.30 &
{[}1.16, 21.44{]} & .029 & .409 & .040 \\
ABAS\_school\_liv\_suppongo & Raven accuracy effectiveness & 13.74 &
{[}0.70, 26.77{]} & .039 & .476 & .095 \\
\end{longtable}
}

\subsection{Navigational Planning Skills in the 5-7 Age
Group}\label{navigational-planning-skills-in-the-5-7-age-group}

To describe baseline navigational planning abilities in the full
analytic sample, PRE Moles Task performance was summarized separately
for 5-, 6-, and 7-year-old children. Means and standard deviations are
reported in the table below. The age effect size is reported as eta2
from one-way ANOVA across the three age groups.

{\def\LTcaptype{none} % do not increment counter
\begin{longtable}[]{@{}
  >{\raggedright\arraybackslash}p{(\linewidth - 10\tabcolsep) * \real{0.1304}}
  >{\raggedleft\arraybackslash}p{(\linewidth - 10\tabcolsep) * \real{0.1739}}
  >{\raggedleft\arraybackslash}p{(\linewidth - 10\tabcolsep) * \real{0.1739}}
  >{\raggedleft\arraybackslash}p{(\linewidth - 10\tabcolsep) * \real{0.1739}}
  >{\raggedleft\arraybackslash}p{(\linewidth - 10\tabcolsep) * \real{0.1739}}
  >{\raggedleft\arraybackslash}p{(\linewidth - 10\tabcolsep) * \real{0.1739}}@{}}
\toprule\noalign{}
\begin{minipage}[b]{\linewidth}\raggedright
Moles Task PRE outcome
\end{minipage} & \begin{minipage}[b]{\linewidth}\raggedleft
Age 5 M (SD)
\end{minipage} & \begin{minipage}[b]{\linewidth}\raggedleft
Age 6 M (SD)
\end{minipage} & \begin{minipage}[b]{\linewidth}\raggedleft
Age 7 M (SD)
\end{minipage} & \begin{minipage}[b]{\linewidth}\raggedleft
Age effect eta2
\end{minipage} & \begin{minipage}[b]{\linewidth}\raggedleft
p raw
\end{minipage} \\
\midrule\noalign{}
\endhead
\bottomrule\noalign{}
\endlastfoot
Working memory span & 1.89 (.60) & 2.14 (1.11) & 2.94 (1.13) & .140 &
\textless{} .001 \\
Working memory trials & 3.00 (.87) & 3.51 (1.39) & 4.54 (1.43) & .155 &
\textless{} .001 \\
Working memory accuracy & 1.22 (1.30) & 1.35 (1.67) & 2.44 (1.49) & .118
& .003 \\
Planning trials & 8.56 (2.55) & 8.30 (.88) & 8.16 (.42) & .015 & .501 \\
Planning accuracy & 13.67 (4.15) & 15.54 (1.26) & 15.42 (1.73) & .073 &
.030 \\
Planning optimal-step trials & 1.22 (1.48) & .38 (1.04) & .52 (1.62) &
.027 & .277 \\
Planning tiles used & 69.89 (20.91) & 74.30 (12.83) & 66.76 (15.13) &
.055 & .071 \\
Planning minimum tiles & 56.22 (16.98) & 60.08 (3.95) & 58.56 (10.82) &
.014 & .520 \\
Planning route efficiency (\%) & 85.05 (10.12) & 86.09 (9.94) & 91.01
(8.46) & .074 & .024 \\
Planning composite & 85.23 (12.49) & 91.38 (5.88) & 93.60 (6.54) & .108
& .005 \\
Working-memory-planning trials & 4.78 (1.56) & 5.05 (1.99) & 5.76 (2.19)
& .035 & .191 \\
Working-memory-planning accuracy & 3.89 (2.32) & 4.57 (2.69) & 5.44
(2.68) & .040 & .148 \\
Working-memory-planning optimal-step trials & 2.56 (.73) & 2.51 (.65) &
2.84 (1.11) & .030 & .248 \\
Working-memory-planning tiles used & 22.67 (12.13) & 24.84 (16.12) &
30.18 (16.77) & .033 & .210 \\
Working-memory-planning minimum tiles & 19.78 (10.80) & 23.59 (13.96) &
26.60 (14.73) & .023 & .335 \\
Working-memory-planning route efficiency (\%) & 87.69 (10.19) & 95.78
(34.27) & 90.13 (9.36) & .018 & .422 \\
Working-memory-planning composite & 56.00 (10.31) & 59.38 (10.40) &
62.04 (10.22) & .034 & .201 \\
\end{longtable}
}

Overall, baseline Moles Task scores showed the expected developmental
pattern: older children tended to achieve higher working-memory scores
and better planning-composite performance, whereas several
working-memory-planning indices showed weaker age separation.

\section{Discussion}\label{discussion}

The final rerun of the paper analyses did not show FDR-corrected
evidence for a generalized chess-training effect across cognitive,
navigational, and timing outcomes. The strongest uncorrected effects
involved forward Corsi span, Moles Task working-memory accuracy, and
Tower of London execution time, but these effects did not survive
correction for the 45 PRE-POST outcomes tested. The age-adjusted ANCOVA
sensitivity analyses led to the same conclusion, suggesting that the
absence of corrected group effects was not explained by baseline age
differences.

At baseline, age was clearly associated with several outcomes.
Seven-year-old children generally performed better than younger children
on Raven accuracy and Moles Task working-memory measures, and they
showed lower Tower of London rule violations. These results support the
developmental sensitivity of the cognitive and navigational measures in
the 5-7 age range.

The predictor analyses suggest that adaptive functioning may be
informative for individual differences in cognitive change. In
particular, teacher-rated ABAS-II Health and Safety was associated with
Tower of London accuracy effectiveness after FDR correction. Other
ABAS-II and BRIEF-2 associations were nominal and should be treated as
exploratory. These findings are consistent with the idea that everyday
adaptive skills may capture individual differences relevant to
structured problem solving, although the present results do not support
strong claims about broad intervention-driven change.

Several limitations should be considered. The sample was relatively
small for detecting modest group x time effects after correction across
many outcomes. The age distribution was unbalanced, with fewer
5-year-old children than 6- and 7-year-old children. In addition, the
full best-available Minefield database was created as an audit resource
but was not yet promoted to the primary paper pipeline; therefore, the
present paper remains based on the final route-efficiency-recalculated
database.

\section{Conclusions}\label{conclusions}

Using the final paper database, the Giant Chess Game protocol was not
associated with FDR-corrected evidence of greater PRE-to-POST
improvement relative to the control group. The analyses did, however,
document robust age-related differences in baseline cognitive and
navigational abilities and identified a corrected association between
teacher-rated adaptive functioning and Tower of London accuracy
effectiveness. Future analyses can use the newly created full
best-available Minefield database as an audit basis before deciding
whether to promote it into the primary paper pipeline.

\section{References}\label{references}

\begin{itemize}
\tightlist
\item
  Insert bibliographic references here.
\end{itemize}




\end{document}
